\section{Language and Compiler Architecture}\label{sec:language}

\Argorix programs group declarations into modules.  Construct families include
agents and capabilities, typed messages and handlers, tools and model calls,
policies, provider contracts, passports, identity and credential metadata,
handshakes, trust ledgers, evidence maps, governance profiles, threat models,
runtime profiles, and release/conformance records.  Some constructs alter
runtime behavior; others are validated descriptions whose execution is
explicitly disabled.

\begin{table}[t]
\centering
\caption{Representative construct families and their boundary in the studied
implementation.}
\label{tab:language-constructs}
\begin{tabularx}{\linewidth}{@{}XXX@{}}
\toprule
Construct family & Representation & Boundary \\
\midrule
Agents and messages & Typed declarations & Implemented \\
Provider contracts & Declarative metadata & Declarative \\
Policies and passports & Semantic plus runtime controls & Implemented \\
ATrust evidence maps & Validated evidence links & Declarative \\
\bottomrule
\end{tabularx}

\end{table}

\Cref{tab:language-constructs} shows why ``represented in the language'' is not
one claim: typed agents and runtime policy have implemented behavior, whereas
ATrust evidence maps and provider contracts remain validated descriptions at
the boundary discussed here.

The pipeline has five relevant stages.  Parsing produces a source-faithful
abstract syntax tree.  Name and module resolution bind declarations without
discarding provenance.  Semantic analysis checks uniqueness, references,
cross-construct constraints, and denied combinations.  Lowering converts only
validated forms to intermediate representation and serializable bytecode.
Finally, the VM loads bytecode, installs explicitly executable providers, and
evaluates policy before scheduling work.

This order matters.  A runtime check alone discovers invalid configuration
after deployment; an unchecked lowering pass can erase the distinction between
an omitted field and an approved field.  Semantic validation rejects, among
other cases, unknown references, duplicated names, external execution without
required approval, inconsistent provider/capability bindings, permissive
runtime profiles, and trust declarations that assert prohibited security
claims.  The conformance corpus contains positive and negative examples for
these boundaries.

\begin{figure}[t]
\centering
\includegraphics[width=\linewidth]{request-sequence.pdf}
\caption{End-to-end request sequence. Semantic validation and policy evaluation
precede operation planning; evidence is produced from the resulting runtime
state.}
\label{fig:request-sequence}
\end{figure}

\Cref{fig:request-sequence} answers RQ2 at the architectural level.  A request
cannot directly select an arbitrary provider.  The VM resolves the agent and
handler, evaluates policy, checks capabilities, and consults the executable
provider registry.  Declarative contracts such as \code{OpenAIProvider} remain
descriptions unless an independently executable adapter is registered.  In the
observed configuration, the registry contains \code{simulated}; attempts to
cross the external boundary are recorded as blocked.

The language also preserves negative statements.  Fields such as
\code{network_allowed=false}, denied key material, mandatory audit, and
fail-closed policy are not comments; their combinations are checked.  However,
the presence of these fields proves only that the compiler accepted the
declaration and the runtime recorded its configured handling.  It does not
prove operating-system containment.
